\section{Structural, procedural and outcome questions}
Structural equivalence concerns strategic opportunities; procedural properties
concern the assignment and selection rules; outcome properties concern realized
allocations. The preceding propositions characterize particular operational
tests, not all three questions simultaneously.

For example, uniformly assigning $n$ otherwise identical participants to $n$
roles gives every participant expected role value $n^{-1}\sum_i m_i$, because
each occupies each role with probability $1/n$. This holds even when the roles
are unequal and each participant strictly prefers one role to another.
Conversely, symmetry of a game does not compel a selection rule to respect it.
An equivariant rule sends a symmetric input to an invariant output or output set;
a deterministic selected point requires a fixed point.

For a finite exact automorphism group $G$, a $G$-invariant profile law gives
equal expected utilities to players in one orbit, by change of summation
indices. Averaging a point's orbit supplies an invariant lottery, but its
independent marginals describe a different procedure. In binary
anti-coordination, the equal lottery on pure equilibria $01,10$ pays each
player 1; the product of its uniform marginals pays each player $1/2$.

The inspected VNM condition is also sensitive to economically meaningful
fixed transfers: singleton-action utilities $(0,0)$ have equal singleton
coalition values, while $(1,-1)$ do not. Both have identical within-role
preferences and transitive cardinal groups. A bridge between coalition-value
fairness and preference invariance therefore needs additional transfer data.
